% ============================================================
%  Harald Høffding, "On the Task and Procedure of the Philosophy of Religion"
%  »Om Religionsfilosofiens Opgave og Fremgangsmaade«
%  Oversigt over det Kgl. Danske Videnskabernes Selskabs Forhandlinger
%  1900, Nr. 5, pp. 411--421.
%  Communicated at the meeting of 2 November 1900.
%  TRANSLATION (English)
%  Translated from the Danish transcription (transcription.tex), which was
%  made from the Royal Danish Academy's digitisation of the 1900 printing.
%  Structure mirrors transcription.tex 1:1, including the page markers, so the
%  two files can be read side by side.
% ============================================================
%
%  TERMINOLOGY. Fixed against this repository's existing Høffding
%  translations, above all »Religionsfilosofi« (1901), which is this paper
%  worked out at length one year later.
%    Værdiens Bestaaen         -> "the persistence of value"
%        The repo's settled term (21 occurrences in religionsfilosofi,
%        4 in relation-som-kategori). B. E. Meyer's 1906 translation of
%        Religionsfilosofi renders it "the conservation of value", and much
%        of the English secondary literature follows Meyer; "persistence" is
%        kept here for consistency within this edition. The two name the same
%        doctrine. "Bestaaen" is likewise "persistence" where it is predicated
%        of anything else (religion's persistence, the persistence of a
%        spiritual tendency), so that the axiom and its applications read as
%        one word in English as they do in the Danish.
%    det religiøse Axiom       -> "the religious axiom" (as in religionsfilosofi)
%    religiøse Forestillinger  -> "religious ideas" (as in religionsfilosofi;
%        note that menneskelige-tanke uses "religious representations" for the
%        same Danish word -- the 1901 sister book governs here)
%    erkendelsesteoretisk      -> "epistemological"
%    aandelig                  -> "spiritual" ("spiritual culture", "the
%        spiritual domain", "the life of the spirit"), following the repo
%    Bevidsthedsliv            -> "the life of consciousness"
%    Sjæleliv                  -> "the life of the soul"
%    Oplevelser                -> "experiences" (Erfaringer likewise; Høffding
%        uses both here and the distinction he draws is carried by the
%        sentence, not by a lexical split)
%
%  QUOTATION MARKS. The printed Danish contains none, anywhere in the item --
%    the one quoted saying, Heraclitus on p. 420, is set unmarked after a
%    colon. That is reproduced here rather than normalised to English quoting
%    convention, so that the translation mirrors the transcription. A reader
%    who expects quotation marks around the saying should know they are absent
%    by the printer's practice, not by an oversight.
%
%  EMPHASIS. There is none in the printed original -- verified by
%    per-character advance-width measurement over all eleven pages. This file
%    therefore contains no \emph{} and no \textit{}, and any that appears
%    later is a defect. See the transcription's header.
%
%  NO GREEK in this item, so textalpha is not loaded; no lists, so enumitem is
%    not loaded. Both are on the playbook's standing checklist and are omitted
%    deliberately, not by oversight.
%
%  THE FOUR PRINTER'S ERRORS. The transcription keeps them as printed and logs
%    each at its site. A translation cannot do the same without turning a
%    compositor's slip into a mistranslation, so the SENSE is translated here
%    and each site carries a "% sic in the Danish:" comment naming the printed
%    form. The four:
%      p. 415  "Forestilliger" for "Forestillinger" -- no effect in English.
%      p. 416  a dropped "som": "det højeste værdifulde, dets Oplevelser have
%              ført det til at kende" -- translated with the relative pronoun
%              the sense requires ("which his experiences have led him to
%              know").
%      p. 419  "af hvilken" for "af hvilket" -- no effect in English, the
%              relative being uninflected.
%      p. 420  "dem gamle Heraklit" for "den gamle Heraklit" -- translated
%              "old Heraclitus".
%
%  TITLE. "Fremgangsmaade" is rendered "Procedure" rather than "Method"
%    because Høffding uses "Metode" as a separate word inside the paper ("Dens
%    Metode er væsentligt analytisk-genetisk"), and the distinction is worth
%    keeping in English.
% ============================================================
\documentclass[12pt,a4paper]{article}
\usepackage[utf8]{inputenc}
\usepackage[T1]{fontenc}
\usepackage{libertinus}
\usepackage{libertinust1math}
\usepackage[english]{babel}
\usepackage[protrusion=true,expansion=true]{microtype}
\usepackage{setspace}
\usepackage[margin=1.2in]{geometry}
\usepackage{fancyhdr}
\usepackage{hyperref}

\hypersetup{
  pdftitle={On the Task and Procedure of the Philosophy of Religion},
  pdfauthor={Harald Høffding},
  hidelinks
}

\pagestyle{fancy}
\fancyhf{}
\fancyhead[LE,RO]{\thepage}
\fancyhead[RE]{\textit{Harald Høffding}}
\fancyhead[LO]{\textit{On the Task and Procedure of the Philosophy of Religion}}
\renewcommand{\headrulewidth}{0pt}

% The four section numbers stand alone on the page in the original, centred,
% with white space above and below and no heading text. Mirrored here.
\newcommand{\snum}[1]{\bigskip\begin{center}#1\end{center}\nopagebreak\medskip}

\setstretch{1.15}

\begin{document}

% Title block mirroring p. 411 of the original. The masthead is given in
% English translation with the Danish original beneath it, since the journal's
% name is how the item is cited.
\begin{center}
  {\footnotesize\scshape Proceedings of the Royal Danish Academy of Sciences\\
   and Letters, 1900.\ No.\ 5.}

  \rule{4em}{0.4pt}

  \bigskip
  {\large\bfseries On the Task and Procedure of the Philosophy of Religion.}

  \medskip
  By

  \smallskip
  {\bfseries Harald Høffding.}

  \smallskip
  {\small (Communicated at the meeting of 2 November 1900.)}

  \medskip
  \rule{4em}{0.4pt}
\end{center}

\medskip

% --- p. 411 ---
\snum{I.}

By philosophy of religion one may understand two different endeavours. One may
mean a philosophizing, a thinking, which springs from the religious
consciousness as an attempt to gain clarity about itself, about the direction
and the object of its innermost need, and about the manner in which that need
is satisfied. Such reflection arises when religion no longer has the character
of immediate feeling or instinct, or of secure repose in tradition; but it
still moves within religion's own presuppositions. Religion is ground and point
of departure, not properly object. I should rather call such an endeavour
religious thinking than philosophy of religion. Some of the greatest minds who
have lived have contributed to this kind of thinking: Augustine, the mystics of
the Middle Ages, Pascal, Jacobi, Kierkegaard. --- By philosophy of religion one
may also understand, and I here understand, a philosophizing, a thinking about
religion as a given psychological and historical phenomenon. Here religion is
object, not subject. That religious thinking itself falls under what is here
the object, since in many ways it sets religion in a clearer light, draws
consequences
% --- p. 412 ---
and opens insights which the unreflective religious consciousness does not
afford.

A philosophy of religion in the sense in which I take the word presupposes that
a religious problem exists. And this presupposes in turn that religion is not
the only form of spiritual culture, and further, that several different
religious standpoints assert themselves in greater or lesser opposition to one
another. The philosophy of religion is comparative: by comparing religion with
other spiritual endeavours, and by comparing different religious tendencies
with one another, it seeks to find religion's proper essence and principle ---
if possible, to formulate this in a single proposition which the various
religions and religious standpoints seek to maintain, in more or less
consistent and pronounced form and in different garbs. Such a proposition must
of course have a certain abstract form, and in this form it exists only in the
philosophy of religion, just as, for example, the laws of the association of
ideas exist only in psychology, while it is nevertheless in accordance with
them that the ideas of practical consciousness combine.

A religious problem cannot arise at any and every time. It presupposes a
differentiation, or a division of labour, in the spiritual domain. It does not
exist in religion's classical ages --- neither in the great ages of beginning,
where religious interest immediately and completely lays claim to all energy,
nor in the great ages of organization, where other spiritual endeavours do
indeed assert themselves, but where they all subordinate themselves to religion
and derive from it their ultimate criteria and regulative principles. In
religion's classical ages there may be religious thinking, but no philosophy of
religion proper. The religious problem presupposes that the unity and the
harmony of the spiritual domain has been broken, or at any rate challenged.
Science and art, morality and social life develop and organize themselves
according to their own laws, on their own independent foundations, and stand as
independent spiritual powers over against religion. And within religion
% --- p. 413 ---
different tendencies likewise assert themselves; indeed the differences may
even concern what is held, from the one side or the other, to be the essential
thing in religion.

The question then becomes what significance religion can have under these
conditions. If it is religion's essence to be the only form of spiritual
culture, and if it lies in religion's essence to appear under one single
determinate form, then it is clear that its time is past once spiritual
division of labour and differentiation have appeared. And in fact such an
assertion is often heard --- while on the other side it is maintained that the
whole intellectual, aesthetic, moral and social development that has taken
place has not altered religion's position within the life of the spirit in the
least. That such opposed assertions can be advanced shows that there really is
a problem here.

It is not likely that this problem will admit of solution. In the spiritual
domain we are far less able than in the physical to distinguish the individual
elements and forces clearly and to find the laws of their interaction, so as to
be able to decide plainly whether the conditions of life for the continued
persistence of a spiritual tendency will remain present or not.

For the investigation of this question it will be of no use to set philosophy,
as a finished system, over against religion. It is philosophy as organ, as
many-sided discussion of the life of the spirit and its conditions, that is
needed here. It may of course have its interest to investigate how a fully
developed philosophical world-view stands to religion. But in the first place
such a system has more the character of art than of science; it stands to
investigative philosophy as the writing of history stands to historical
criticism. And in the second place a philosophical system will have value only
if, already in its coming-to-be and its preparation, it has taken due account
of such a spiritual fact as religion is.
% --- p. 414 ---
It will then be natural, in discussing the religious problem, to employ
philosophy in its investigative and preparatory shape. Philosophy in this sense
is a critical and comparative discipline, which seeks out the fundamental forms
and fundamental presuppositions of the life of consciousness, and the
conditions of their persistence and their validity. Its method is essentially
analytic-genetic, in that from the finished works and forms it goes back to the
forces and elements out of which they have arisen. An investigation of religion
of this kind may suitably be laid out so as to consist of an epistemological, a
psychological and an ethical part. --- I shall attempt to give a condensed
survey of such an investigation.

\snum{II.}

In religion's classical ages, in which it was all in all, religious ideas
furnished a satisfying explanation and understanding of existence. There has
now developed a need for understanding which is not met by these ideas. The
individual phenomenon, the individual event in the spiritual and in the
material world, is understood (on this new conception of understanding) by
being brought into the closest continuous connection with other phenomena and
events, each of which can be exhibited in experience just as well as those
whose explanation was asked for. It is required that the cause, no less than
the effect, shall be an experience. It is now acknowledged more and more
(though this has cost hard and repeated struggles) that it is not religion's
business to furnish an understanding of this kind. And when the two modes of
explanation collide over one and the same phenomenon, the question becomes how
they are to be reconciled, and whether on the whole they can be. This much,
however, seems on the whole to be generally conceded: that religious ideas do
not have their significance and their value in being required to afford full
satisfaction of the need for knowledge. It
% --- p. 415 ---
is then probable that religious ideas owe their origin and their persistence to
sides of the life of consciousness other than the intellectual.

One might suppose that it was only the individual, particular phenomena that
science could explain, while the concluding questions, the ultimate riddles,
had to be sought answered in religion. But if it now turns out, on closer
critical investigation, that religion is as little able as science to solve the
ultimate riddles, to give us really conclusive and consistent concepts, then
the question of the significance of religious ideas returns with renewed force.
It turns out (as it would be too lengthy to establish here in detail) that
religious ideas are essentially due to analogy, without the analogies that here
lie at the basis being capable of being formed and applied in so clear and
thoroughgoing a manner as the analogies of which scientific thinking makes use
in various domains. This is no accident; for as it stands with religious
% sic in the Danish: the page prints "Forestilliger" for "Forestillinger",
% the n dropped. No effect in English.
ideas, so it stands also with the thoughts by which speculative philosophy, in
the various metaphysical systems, has attempted to give a conclusion to the
understanding of existence. There comes a point here at which thought is
relieved by poetry. Religious ideas do indeed have throughout a figurative
character. They are great symbols in which spiritual states of another kind and
compass than the purely intellectual have sought and found expression. --- At
this point the epistemological philosophy of religion leads naturally over into
the psychological.

\snum{III.}

This part of the science of religion will more and more become the most
important, and it is the part that has offered me the greatest interest. Here
we are led away from the strife between dogmatics and criticism, back to the
life of the soul in its inmost and most
% --- p. 416 ---
% sic in the Danish: the page prints "det højeste værdifulde, dets Oplevelser
% have ført det til at kende" -- a dropped "som" before "dets". Translated
% here with the relative pronoun the sense requires.
concentrated forms. Religion cannot be made or constructed. It grows out of
life, springs from man's fundamental mood in the struggle of life, from his
will to hold fast to what is of highest value, which his experiences have led
him to know. We do not always come to know it so well from this side with the
help of the history of religion, significant as that of course otherwise is. In
the history of religion one does not so easily get to see the forces of the
soul in their original working. It occupies itself more with religion's works
than with religion itself --- more with the outward occasions and conditions
than with the inner life --- more with the great types, the forms common to
large groups of men, than with the individual processes. The study of
individual religiosity, as this can be known from the biographies, and
especially the autobiographies, of religious personalities, is therefore of the
greatest importance for the psychology of religion. Augustine's Confessions,
and the autobiographies of Suso and of St.\ Teresa, are psychological
source-documents of the first rank. There is of course, and particularly in the
religious domain, a close connection between individual and society, and an
investigation of the relation between individual religiosity and religious
tradition is therefore one of the most important chapters in the psychology of
religion. Individual religiosity belongs moreover to the history of religion as
well, although that discipline has hitherto occupied itself almost exclusively
with the great types and has disregarded the manifold ways in which they are
shaded in single individuals. And yet it is really only through these shadings
that the operative forces and elements come clearly to light.

There are two concepts which it is above all the business of the psychology of
religion to determine: religious experience and religious faith.

The peculiarity of religious experience I find in this, that whereas other
experiences concern either the given reality --- which holds of the experiences
we make through the various senses --- or the value this reality has for us ---
which holds of the experiences due to organic vital feeling and to aesthetic,
intel\-%
% --- p. 417 ---
lectual and ethical feeling --- religious experience, or religious feeling,
concerns the relation between value and reality. It is determined by the fate
of what is of value in the world of reality. --- The peculiarity of religious
faith I find in this, that amid the course of actual events it holds fast to
the persistence of value. It is a deep trust, a continuity in the mind through
all fluctuations. Its objective content is a continuity too: its fundamental
assumption is --- anthropomorphically expressed --- a faithfulness in
existence, which shows itself in the fact that the values persist, however much
their particular forms may change.

The correctness of this psychological characterization may be established along
several routes. In the first place, it is possible, on the basis of the marks
given, to describe a direction of feeling and of will which must naturally
develop under human conditions of life and in virtue of the general
psychological laws. It is a psychologically possible state. In the second
place, by employing this description, and calling what it describes religion,
one will be able to exhibit clearly and distinctly the difference of religious
feeling from, and its relation to, the other kinds of feeling with which it is
akin. And in the third place, an investigation of the most important
historically given types of religion shows that the marks named are the most
important and the pervasive ones.

The proposition of the persistence of value may be said to be the religious
axiom. It expresses a fundamental thought which the religious consciousness
holds fast under the most various forms, and which, in the debate between
different religious standpoints, emerges more or less distinctly as ideal and
as standard. The difference between religious standpoints rests partly on the
values in whose persistence one believes, partly on the clarity with which the
thought of the persistence of values emerges and is held fast, partly on the
different ways in which it is expressed. For the single determinate religion
the essential thing is no doubt the determinate form in which it expresses the
axiom. But for the comparative
% --- p. 418 ---
investigation the point is to find the typical, the theme, which in the
particular case may be varied in many ways.

With the help of the hypothesis thus grounded it now becomes possible to
discuss the question that is the kernel of the religious problem, namely that
of religion's persistence under the altered conditions which the division of
labour in the spiritual domain brings with it. If the hypothesis is right, a
distinction will have to be drawn between religion's form and its kernel. The
forms --- myths and dogmas --- may change and the kernel yet persist. That is
why purely intellectual criticism is in the end of such subordinate importance.
Because certain religious ideas fall away, the essential motive for the
formation of those ideas does not fall away. It may express itself in new
effects, just as the need to find causes does not fall away because certain
causal explanations have proved invalid. Religion's time would be past only if
it could be proved impossible to hold fast to the assumption of the persistence
of value.

The religious axiom cannot be derived from science. Scientific thinking has
completed its work when it has exhibited as many relations of identity, of
rationality and of causality as are possible for it at a given time. The picture
of the world that thereby arises retains its validity, whatever the fate of the
values in the world of reality may turn out to be. But the religious axiom, on
the other hand, need not lead to conflict with the scientific picture of the
world, although in fact this has often been the case. The possibility is not
excluded that it might be expressed in a way that excluded such a conflict. Or
the very forms and laws of the actual picture of the world might stand
precisely as the ways and the means that make the persistence of value
possible. It would be the ideal if the standpoint of value and the standpoint
of reality could in the end coincide. From the standpoint of critical
philosophy the question cannot be held back, whether the difference we draw,
and must draw, between value and reality is not merely subjective and human.
This
% --- p. 419 ---
difference might well rest on the fact that both our concept of value and our
concept of reality are imperfect, unconcluded and inconcludable. We can derive
neither value from reality nor reality from value. Our thought is then directed
to the possibility of a principle from which the persistence of value and the
% sic in the Danish: the page prints "af et Princip, af hvilken" -- the neuter
% "Princip" requires "hvilket". No effect in English, the relative pronoun
% being uninflected.
fixed connection of reality might spring at one and the same time. But it is a
thought that cannot be carried through scientifically.

There will always be matter enough for strife in the religious domain. There
may be strife about which values they are in whose persistence one is to
believe, and there will be strife about the right forms of the faith in the
persistence of those values. But the fundamental religious point of contention
will be whether the religious axiom itself can be maintained. More than a
possibility of this cannot be established theoretically, if it is right that
all faith is a matter of life. But life is perhaps richer and more fruitful
than narrow-minded doubts and narrow-minded dogmas suppose.

\snum{IV.}

The ethical philosophy of religion naturally forms the conclusion of the whole
investigation. It examines two main questions, namely how far religion is the
foundation of ethics, and what significance religion has --- even should that
not be the case --- as a form of spiritual culture. Both questions may gain in
clarity with the help of the hypothesis that faith in the persistence of value
is the essential thing in religion.

The values in whose persistence a man believes he must first know for himself
from experience. When ethical ideas appear in religious form --- when, for
example, the gods appear with ethical attributes --- the presupposition is that
these ideas or attributes have proved their value in human life itself. Only so
can a man connect any meaning with them when they are used to express the
divine. Historians of religion therefore lay great weight on the development of
men's ethical
% --- p. 420 ---
ideas as the condition of the most important opposition within the history of
religion, the opposition between nature-religion and ethical religion. Both
historically and purely logically considered, the ethical has therefore
priority in relation to the religious. And this emerges particularly in the
fact that religious debates lead back in the end to ethical oppositions, to a
difference of valuation --- that is, to a difference in respect of the values
in whose persistence one believes.

But even if this is right, the relation between religion and ethics is not
thereby exhausted. From a purely ethical standpoint it must still be asked:
what value is there in believing in the persistence of value? That is the
question of religion's significance as a form of spiritual culture. The
treatment of this question calls for a thorough psychological and sociological
investigation. Psychologically considered, such a faith has its significance in
strengthening the will and in the widening of the horizon to which it leads.
Faith in the persistence of value allows us --- even where we cannot follow the
metamorphoses of value in detail --- to view existence as a great system of
potential values, which in part are to be converted through human labour into
actual values. Through its inmost root the ethical consciousness here hangs
together with the religious. Faith in the persistence of value acquires in
ethics a significance similar to that which ideas, anticipations and hypotheses
have in the theoretical domain. An anxious clinging to given experiences easily
shuts one out from making new experiences. Already old Heraclitus said: If you
% sic in the Danish: the page prints "Allerede dem gamle Heraklit" -- "dem"
% for "den". The saying that follows is set without quotation marks in the
% original, after the colon; that is reproduced here. (Heraclitus, DK B18.)
do not hope, you will not find out that which could not be hoped for.
Sociologically, religion has its great significance through its organization
and through the philanthropic activity it calls forth.

A more or less sharp relation of opposition between religion and ethics arises
not only when religion sets up commands or prohibitions which ethically cannot
be acknowledged. The most characteristic opposition arises when faith in the
persistence of value itself becomes a hindrance to the work of
% --- p. 421 ---
discovering and producing values, by laying claim to energy and time that might
otherwise have been used in the service of that work. If there are those whose
whole energy is spent in such work, the question becomes whether it diminishes
the value of their personality that they do not --- at any rate not with full
consciousness --- believe in the persistence of value in any form whatever.
This is perhaps the most decisive form in which the religious problem can
present itself.

It seems, then, that the view in the philosophy of religion which has here been
indicated in broad outline has not only psychological and historical grounding,
but is also able to furnish important points of view for the struggle in the
religious domain.
% End of the article, mirroring the Danish. In the original the text stops
% about a third of the way down p. 421, followed by a short centred decorative
% double rule, which is not set here either. No date line, no signature.

\end{document}
